\subsection{A logical roadmap} There are several recent results in the literature which inform the Toda and Whitney
presentations we prove in this paper.
Since we do not attempt to give self-contained
proofs, we provide next a roadmap of the logical
implications we rely on, which a reader may find useful.
 %(Each of these will be recalled again during the relevant proofs.)

Our proofs of both the Toda presentation, and the Whitney presentation,
for $\QK_T(\Fl(\bfr,n))$, from Theorem~\ref{thm:main} (resp.\ \eqref{E:QKWhitney}, see Remark~\ref{rmk:QKW})
rely on the following: the Toda (resp.\ Whitney) presentation for $\QK_T(\Fl(n))$;
Kato's push forward homomorphism from~\eqref{E:kato-push}; and the key
technical result from Proposition~\ref{prop:MNS}, proved in~\cite{maeno.naito.sagaki:QKideal},
which in turn relies on Kato's work~\cite{kato:loop}.

There are two proofs of the Toda presentation for $\QK_T(\Fl(n))$, one in
\cite{maeno.naito.sagaki:QKideal}, relying on~\cite{kato:loop}, and another
which may be deduced from~\cite{act:2017}, relying on results from
\cite{givental_lee} and~\cite{Iritani:2013qka};~cf.~Appendix~\ref{sec:toda}.

There are also two proofs of the Whitney presentation for $\QK_T(\Fl(n))$.
One was recently announced by Huq-Kuruvilla
\cite{huq2024quantum} (for all rings $\QK_T(\Fl(\bfr,n))$), and
it uses the technique of abelian-nonabelian correspondence, independent
of Kato's results. Another proof of the Whitney presentation for
 $\QK_T(\Fl(n))$ is given in~\cite[Section~6]{GMSXZZ:QKW}.
It relies on the recent proof of the quantum~K divisor axiom~\cite{LNSX:QKdiv},
and ultimately on Kato's results.

\section{Preliminaries}
\subsection{Equivariant K theory of Grassman bundles}
Let $T$ be a linear algebraic group. For any projective $T$-variety $Z$,
let $\K_T(Z)$ be the equivariant $\K$-theory
ring, defined as the Grothendieck ring of $T$-equivariant algebraic vector
bundles. This ring is an algebra over $\K_T(\pt)$, the representation ring of
$T$. Let $\euler{Z}\colon \K_T(Z) \to \K_T(\pt)$ be the pushforward map along the
structure morphism.

For $E \to Z$ a $T$-equivariant
vector bundle of rank $\operatorname{rk}E$, we denote by
\[ \lambda_y(E) := 1 + y [E] + \dots + y^{\operatorname{rk}E} \bigl[\wedge^{\operatorname{rk}E} E\bigr] \in \K_T(Z)[y] \]
the Hirzerbruch $\lambda_y$ class of $E$. This class is multiplicative for short exact sequences. {In an abuse of notation, we often write $E$ for the class $[E]$ in $\K_T(Z)$.} Note that for a rank $e$ equivariant vector bundle $E$, and a character ${\rm e}^\chi \in \K_T(\pt)$,
\[ \lambda_y( {\rm e}^\chi \otimes E) = \lambda_{y {\rm e}^{\chi}}(E) = \sum_{i=0}^{e} y^i {\rm e}^{i \chi} \otimes \wedge^i E . \]
As is customary, we will often remove the $\otimes$ symbol from the notation.

Denote by
$\pi\colon \mathbb{G}(r,E) \to Z$ the Grassmann bundle over $Z$.
It is equipped with a tautological sequence
$0 \to \underline{\cS} \to \pi^* E \to \underline{\cQ} \to 0$ over
$\mathbb{G}(r,E)$. The following result follows
from~\cite[Proposition~2.2]{kapranov:Gr}, see also~\cite[Proposition~3.2 and Corollary~3.3]{gu2022quantum}.
(Kapranov proved this when $Z=\pt$; the relative version follows immediately
using that $\pi$ is a $T$-equivariant locally trivial fibration.) We only state the special cases that will be used in this paper. See the above references for the full generality.

\begin{Proposition}[Kapranov]\label{prop:relativeBWB}
 There are the following isomorphisms of $T$-equivariant vector bundles:
 \begin{enumerate}\itemsep=0pt
 \item[$(1)$] for all $i \ge 0$, $\ell>0$ the higher direct images, $R^i \pi_* \bigl(\wedge^\ell\underline{\cS}\bigr) = 0$;
 \item[$(2)$] for all $\ell\geq 0$, \[ R^i \pi_*\bigl(\wedge^\ell\underline{\cQ}\bigr) = \begin{cases} \wedge^\ell E, & i=0, \\ 0, & i>0 . \end{cases}\]
 \end{enumerate}
 \end{Proposition}

\subsection{(Equivariant) quantum K theory of flag varieties}\label{sec:prelim}
Let $\mathbf{r}= (r_1, \dots, r_k)$. We consider \[X=\Fl(\bfr,n),\] which parametrizes flags of vector spaces
$F_{1} \subset F_{2} \subset \dots \subset F_{k} \subset \C^n$
with $\dim F_{i} = r_i$ for~${1 \le i \le k}$.

Let $M_{d,n} := \Mb_{0,n}(X,d)$ be the moduli space of genus zero degree
$d$ stable maps to $X$ with~$n$ marked points. Given classes
$a_1, \dots, a_n \in \K_T(X)$,
define the $\K$-theoretic Gromov--Witten invariants by
\[ \langle a_1, \dots, a_n \rangle_{d} =
\euler{M_{d,n}}\Biggl(\prod_{i = 1}^n\ev_i^*(a_i)\Biggr)\in \K_T(\pt) . \]
Non-equivariant Gromov--Witten invariants are obtained by replacing $T$ with the
trivial group; these Gromov--Witten invariants are integers.

For $d=(d_1,\dots,d_k)\in\HH_2(X,\Z)\cong\Z^k$, we write $Q^d$ for \smash{$\prod_{i=1}^k Q_i^{d_i}$}. Here $Q_i$ corresponds to the Poincar{\'e} dual of the first Chern class $-c_1(\det \cS_i)$. Following~\cite{givental:onwdvv,lee:QK}, the $T$-equivariant (small) quantum
K theory ring is
\begin{equation*}%\label{eqn:QKTfree}
\QK_T(X) = \K_T(X) \otimes_{\K_T(\pt)} \K_T(\pt)[\![Q]\!]
\end{equation*}
as a $\K_T(\pt)[\![Q]\!]$-module.
It is equipped with a commutative, associative
product, denoted by $\star$, which is
determined by the condition
\begin{equation*}%\label{eqn:def}
 (\!(\sigma_1\star\sigma_2,\sigma_3)\!)=\sum_d
 Q^d \gw{\sigma_1,\sigma_2,\sigma_3}{d}\qquad\text{for all $\sigma_1,\sigma_2,\sigma_3\in \K_T(X)$},
\end{equation*}
where
\[
 (\!(\sigma_1,\sigma_2)\!)\coloneqq \sum_{d}Q^d \gw{\sigma_1,\sigma_2}{d}
\]
is the quantum $\K$-metric.

It was proved in~\cite{anderson2022finite, kato:loop} that for $\sigma_1,\sigma_2\in\K_T(X)$, the product $\sigma_1\star\sigma_2$ can always be expressed as a polynomial in $Q$ with coefficients in $\K_T(X)$. It follows that
\[
 \QK_T^\poly(X)\coloneqq\K_T(X) \otimes_{\K_T(\pt)}\K_T(\pt)[Q]
\]
is a subring of $\QK_T(X)$.

Let $Y=\Fl\bigl(r_1,\dots, \widehat{r}_i,\dots, r_k;n\bigr)$ and $\pi\colon X\to Y$ be the natural map. Let also $\widehat{\mathbf{r}} = \bigl(r_1, \dots,\allowbreak \widehat{r}_i,\dots, r_k\bigr)$. The following theorem is a specialization of results proved in~\cite{kato2019quantum}.

\begin{Theorem}[Kato]\label{thm:kato-morphism}
 There is a surjective ring homomorphism
 \[
 \Phi\colon\ \QK^\poly_T(X)\to\QK^\poly_T(Y)
 \]
 given by $\sigma\mapsto{\pi}_*\sigma$ for all $\sigma\in\K_T(X)$ and
 \[
 Q_j\mapsto\begin{cases}
 Q_j, & j\neq i,\\
 1, & j=i
 \end{cases} \qquad \text{for $1\leq i\leq k$}.\]
\end{Theorem}
It follows from Theorem~\ref{thm:kato-morphism} that
Kato's homomorphism extends naturally to
\begin{equation}\label{E:Kato-loc-hom}
\Phi\colon \QK_T^{\loc(\mathbf{\hat r})}(X) \to
\QK_T^{\loc(\mathbf{\hat r})}(Y) ,
\end{equation}
where $\loc(\mathbf{\hat r})$ indicates localization at the multiplicative set generated by
$1- Q_j$ for $j \neq i$.

\subsection[The Toda presentation for Fl(n)]{The Toda presentation for $\boldsymbol{\Fl(n)}$}

The variety $\Fl(n)=\Fl(1,\dots,n-1;n)$ is equipped with tautological vector bundles
\[
 0=\cS_{0}\subset\cS_{1}\subset\dots \subset \cS_{n-1}\subset \cS_{n}=\C^n ,
\]
where $\cS_j$ has rank $r_j$. It can also be viewed as $\SL_n/B$, where $B\subset \SL_n$ is a Borel subgroup. Let~${T\subseteq B}$ be a maximal torus in $\SL_n$.

The following is the main result of~\cite{maeno.naito.sagaki:QKideal} (see Remark~\ref{rmk:mns} for more details). The relation~\eqref{eqn:toda} can also be recovered from the connection between the $J$-function of the full flag variety and the relativistic Toda lattice established by Givental and Lee in~\cite{givental_lee}. This observation was made in the unpublished note~\cite{act:2017} of Anderson--Chen--Tseng, but removed from the published version of their paper. For the sake of completeness, we give a brief account in Appendix~\ref{sec:toda}.

\begin{Theorem}\label{thm:toda}
The ring $\QK_T(\Fl(n))$ is isomorphic to $R'[\![Q]\!]/J'_Q$, where $R'$ is equal to the ring \smash{$\K_T(\pt)\bigl[P_1^{\pm },\dots,P_{n}^\pm\bigr]$} and the ideal $J'_Q\subset R'\dbb{Q}=R'\dbb{Q_1,\dots,Q_{n-1}}$ is generated by the coefficients $y$ in
 \begin{gather}\label{eqn:toda}
 \lambda_y(\C^n) -
 \begin{vmatrix}
 1+y\frac{P_1}{P_0} & y\frac{P_2}{P_1}Q_1 & & & & \\[0.9mm]
 1 & 1+y\frac{P_2}{P_1} & y\frac{P_3}{P_2}Q_2 & & & \\[0.9mm]
 & 1 & 1+y\frac{P_3}{P_2} & y\frac{P_{4}}{P_{3}}Q_3 & & \\
 & & \ddots & \ddots & \ddots & \\
& & & 1 & 1+y\frac{P_{n-1}}{P_{n-2}} & y\frac{P_n}{P_{n-1}}Q_{n-1}\\[1mm]
& & & & 1 & 1+y\frac{P_n}{P_{n-1}}
\end{vmatrix}^\star,
\end{gather}
here $P_0=1$ by convention, and $\lambda_y(\C^n)\in\K_T(\pt)[y]$.

More precisely,
there exists a $\K_T(\pt)[\![Q]\!]$-algebra isomorphism $\Psi'\colon R'[\![Q]\!]/J'_Q \to \QK_T(\Fl(n))$ that sends $P_j$ to $\det \cS_j$ for all $j=1,\dots,n$.
\end{Theorem}

\begin{Remark} \label{rmk:mns}
 Theorem~\ref{thm:toda} is proved in~\cite{maeno.naito.sagaki:QKideal} using results of Kato~\cite{kato:loop} based on the semi-infinite
 flag variety.
 		The connection between our statement of Theorem~\ref{thm:toda} and that of~\cite{maeno.naito.sagaki:QKideal} is seen as follows. Define the \textit{Toda polynomials} \smash{$T^{(n)}_k$} for $k=1,\dots,n$ by
		\begin{align*}
		T^{(n)}_k = \sum_{0= i_0<\dots<i_k\leq n}\ \prod_{s=1}^{k}\frac{P_{i_s}}{P_{i_{s}-1}}\left(1-Q_{i_s-1}\right)^{1-\delta_{i_s-i_{s-1},1}}.
		\end{align*}
		These elements of $\Z\bigl[P_1^{\pm},\dots,P_{n}^{\pm}\bigr][\![Q]\!]$ (where $P_0=1$ and $Q_0=0$ by convention) are symbols of the finite-difference Toda Hamiltonians~\cite{etingof} (see also~\cite{act:2017,gloI,givental_lee,koroteev}).
		
		Letting \smash{$T^{(n)}=\sum_{k=0}^n T^{(n)}_k y^k$} where \smash{$T^{(n)}_0=1$}, we claim that $T^{(n)}(y)$ is equal to the determinant of the matrix appearing in the Toda relations~\eqref{eqn:toda}, namely,
		\begin{align*}
 T^{(n)}(y)=\begin{vmatrix}
 1+y\frac{P_1}{P_0} & y\frac{P_2}{P_1}Q_1 & & & & \\[0.9mm]
 1 & 1+y\frac{P_2}{P_1} & y\frac{P_3}{P_2}Q_2 & & & \\[0.9mm]
 & 1 & 1+y\frac{P_3}{P_2} & y\frac{P_{4}}{P_{3}}Q_3 & & \\
 & & \ddots & \ddots & \ddots & \\
& & & 1 & 1+y\frac{P_{n-1}}{P_{n-2}} & y\frac{P_n}{P_{n-1}}Q_{n-1}\\[1mm]
& & & & 1 & 1+y\frac{P_n}{P_{n-1}}
\end{vmatrix}^\star.
		\end{align*}
		This is verified by showing that $T^{(n)}$ satisfies the recursion
		\begin{equation*}%\label{E:W-Trec}
T^{(n)}= T^{(n-1)} \biggl(1+ y \frac{P_n}{P_{n-1}}\biggr) - y Q_{n-1} \frac{P_n}{P_{n-1}} T^{(n-2)} ,
		\end{equation*}
		and then applying Lemma~\ref{L:r-det} below (with $n$ playing the role of $j$ there).
		\end{Remark}

 \begin{Remark}
 Upon the specialization $Q_1=\dotsm=Q_{n-1}=0$, the Toda presentation $R'[\![Q]\!]/J'_Q\cong\QK_T(\Fl(n))$ becomes the Borel presentation
 \[
 \K_T(\pt)\bigl[P_1^\pm,\dots,P_n^{\pm}\bigr]/J\cong \K_T(\Fl(n)),
 \]
 where $J$ is the ideal generated by the coefficients of $y$ in
 \[
 \lambda_y(\C^n)-\sideset{}{^\star}\prod_{j=0}^{n-1} (1+yP_{j+1}/P_j)
 \]
 and~$P_j$ corresponds to $\det \cS_j$ for all $j=1,\dots,n$.
 \end{Remark}
 \begin{Lemma}
 \label{L:r-det}
 Suppose $U_j$ for $0\le j\le k+1$ and $A_j$, $B_j$ for $0\le j\le k$ are elements of a~commutative ring with $1$ such that the $U_j$ satisfy the recursion
 \begin{equation*}
 %\label{E:recursion}
 U_{j+1}=A_jU_{j}-B_jU_{j-1},\qquad 0\le j\le k
 \end{equation*}
 with initial conditions $U_0=1$, $U_{-1}=0$. Then, for all $0\le j\le k+1$, one has
 \begin{equation}
 \label{E:r-det}
 U_j=
 \begin{vmatrix}
 A_0 & B_1 & & & \\
 1 & A_1 & B_2 & & \\
 & \ddots & \ddots & \ddots & \\
 & & 1 & A_{j-2} & B_{j-1}\\
 & & & 1 & A_{j-1}
 \end{vmatrix}.
 \end{equation}
 \end{Lemma}

 \begin{proof}
 One simply expands along the last row or column to see that the determinant in~\eqref{E:r-det} satisfies the recursion. Observe that the initial values $U_0=1$ and $U_1=A_0$ agree. This completes the proof.
 \end{proof}

 Before finishing this section, we record the following, which follows from~\cite[Proposition~5.2]{maeno.naito.sagaki:QKideal}.

 \begin{Proposition}[Maeno--Naito--Sagaki]\label{prop:MNS}
 In $\QK_T(\Fl(n))$, the following relations hold:
 \begin{gather*}
 \det \cS_i \star \det \cS_j/\cS_i = (1-Q_i) \det \cS_j, \qquad 1 \le i < j \le n.
 \end{gather*}
 \end{Proposition}

\section[Toda-type presentations for the equivariant quantum K theory of partial flag varieties]{Toda-type presentations for the equivariant quantum\\ K theory of partial flag varieties}%\label{sec:partial}

To begin, we observe that the Toda presentation in Theorem~\ref{thm:toda} can be rewritten as follows.
\begin{Corollary}\label{cor:fln}
 The ring $\QK_T(\Fl(n))$ is isomorphic to $R\dbb{Q}/J_Q$, where
 \[
 R=\K_T(\pt)\bigl[Y^{(0)}, \dots, Y^{(n-1)}\bigr],
 \]
 and $J_Q\subset R\dbb{Q}=R\dbb{Q_1,\dots,Q_{n-1}}$ is generated by the coefficients of $y$ in
 \begin{gather*}%\label{eqn:toda-gen}
 \lambda_y(\C^n) - \\
 \begin{vmatrix}
 1+yY^{(0)}\frac{1}{1-Q_0} & yY^{(1)}\frac{Q_1}{1-Q_1} & & & & \\[1mm]
 1 & 1+yY^{(1)}\frac{1}{1-Q_1} & yY^{(2)}\frac{Q_2}{1-Q_2} & & & \\ % & 1 & 1+yY^{(2)}\frac{1}{1-Q_2} & yY^{(3)}\frac{Q_3}{1-Q_3}Q_3 & & \\
& \ddots & \ddots & \ddots \\ & & 1 & 1+yY^{(n-2)}\frac{1}{1-Q_{n-2}} & yY^{(n-1)}\frac{Q_{n-1}}{1-Q_{n-1}}\\[1mm] & & & 1 & 1+yY^{(n-1)}\frac{1}{1-Q_{n-1}}
\!\!\!\end{vmatrix}^\star
\end{gather*}
with the convention that $Q_0=0$.

More precisely,
% for $\mathbf{r}=(1,2,\dots,n-1)$,
there exists a $\K_T(\pt)[\![Q]\!]$-algebra isomorphism $\Psi\colon R[\![Q]\!]/J_Q \to \QK_T(\Fl(n))$ that sends $Y^{(j)}$ to $\cS_{j+1}/\cS_j$ for $j=1,\dots,n-1$.
\end{Corollary}
\begin{proof}
 Identifying $P_{j+1}/P_j$ with $Y^{(j)}/(1-Q_j)$ gives an isomorphism between $R[\![Q]\!]/J_Q$ and $R'[\![Q]\!]/J'_Q$. More precisely, define a $\K_T(\pt)[\![Q]\!]$ homomorphism $\Phi\colon R'[\![Q]\!]/J'_Q\to R[\![Q]\!]/J_Q$ by
\[
 \Phi(P_j)=\prod_{i=0}^{j-1}\frac{Y^{(i)}}{1-Q_i}, \qquad 1\leq j\leq n-1.
\]
Note that in $R[\![Q]\!]/J_Q$, we have \smash{$\det\C^{n}=\prod_{i=0}^{n-1}Y^{(i)}/(1-Q_i)$}, which implies all $Y^{(j)}$ are invertible. Since the relations match, the homomorphism $\Psi$ is well-defined and injective. Since \smash{$(1-Q_j)P_{j+1}/P_j$} is sent to $Y^{(j)}$ for $0\leq j\leq n-1$, it is also surjective. Finally, the~geometric interpretation follows from Proposition~\ref{prop:MNS}.
\end{proof}

Next, we generalize Corollary~\ref{cor:fln} to all partial flag varieties utilizing Theorem~\ref{thm:kato-morphism}, Proposition~\ref{prop:MNS}, and the Nakayama-type result from~\cite{GMSXZZ:QKW, gu2022quantum}.

\begin{Theorem}\label{thm:whit-short}
 In $\QK_T(\Fl(\bfr,n))[y]$, the following relation hold:
 \begin{align}\label{eqn:whit-short}
 \lambda_y(\C^n)-
 &\begin{vmatrix}
 A_0 & B_1 & & & \\
 1 & A_1 & B_2 & & \\ & \ddots & \ddots & \ddots & \\ & & 1 & A_{k-1} & B_{k}\\ & & & 1 & A_{k}
 \end{vmatrix}^\star,
 \end{align}
 where \[
 B_j=y^{r_{j+1}-r_j}\frac{Q_j}{1-Q_j}\det(\cS_{j+1}/\cS_j),\qquad
 A_j=\lambda_y(\cS_{j+1}/\cS_j)+B_j.
\]
\end{Theorem}
\begin{proof}
Let $X=\Fl(r_1,\dots,r_k;n)$, $Y=\Fl\bigl(r_1,\dots, \widehat{r}_i,\dots, r_k;n\bigr)$, and $\pi\colon X\to Y$ be the natural map. Let
\[
 0=\cS_0\subset\cS_1\subset\dots\subset\cS_k\subset\cS_{k+1}=\C^n
\]
be the sequence of tautological bundles on $X$. Note that all but $\cS_i$ are pulled back from $Y$. With a slight abuse of notation, we denote the sequence of tautological bundles on $Y$ by
\[
 0=\cS_0\subset\cS_1\subset\dots\cS_{i-1}\subset\cS_{i+1}\subset\dots\subset\cS_k\subset\cS_{k+1}=\C^n.
\]
Note that the elements $B_1,\dots,B_{i-2},B_{i+1},\dots,B_k$ as well as $A_1.\dots, A_{i-2},A_{i+1},\dots,A_k$ in the ring $\QK_T(X)[y]$ stay the same under pushforward along $\pi$. By a slight abuse of notation, we~also think of them as elements of $\QK_T(Y)[y]$.

By induction, we assume that relation~\eqref{eqn:whit-short} holds for $X$, i.e.,
\begin{gather}\label{eqn:whitX}
 \lambda_y(\C^n)-\begin{vmatrix}
 A_0 & B_1 & & & & & & & & & \\
 1 & A_1 &B_2 & & & & & & & & \\
 & \ddots & \ddots & \ddots & & & & & & & \\
 & & \ddots & \ddots & B_{i-2} & & &\\ & & & 1 & A_{i-2} & B_{i-1} & & & & &\\ & & & & 1 & A_{i-1} & B_i & & & &\\
 & & & & & 1 & A_i & B_{i+1} & & &\\
 & & & & & & 1 & A_{i+1} & \ddots & &
 \\
 & & & & & & & \ddots & \ddots & \ddots & \\
 & & & & & & & & 1 & A_{k-1} & B_{k}\\
 & & & & & & & & & 1 & A_{k}
 \end{vmatrix}^\star
\end{gather}
holds in \smash{$\QK_T^{\loc(\mathbf{r})}(X)[y]$}
for $1 \le j \le k$, and we will show that the (localized) Kato's pushforward~\eqref{E:Kato-loc-hom} of this relation gives relation~\eqref{eqn:whit-short} on $Y$.

Relation~\eqref{eqn:whit-short} on $Y$ reads
\begin{gather}\label{eqn:whitY}
 \lambda_y(\C^n)-\begin{vmatrix}
 A_0 & B_1 & & & \\
 1 & A_1 & B_2 & & \\ & \ddots & \ddots & \ddots & \\
 & & \ddots & \ddots & B_{i-2}\\ & & & 1 & A_{i-2} & B'_{i-1}\\ & & & & 1 & A'_{i-1} & B_{i+1}\\ & & & & & 1 & A_{i+1} & \ddots
 \\ & & & & & & \ddots & \ddots & \ddots \\ & & & & & & & 1 & A_{k-1} & B_{k}\\ & & & & & & & & 1 & A_{k}
 \end{vmatrix}^\star,
\end{gather}
where
\begin{gather*}
 B_{i-1}'=y^{r_{i+1}-r_{i-1}}\frac{Q_{i-1}}{1-Q_{i-1}}\det\left(\cS_{i+1}/\cS_{i-1}\right),\qquad A_{i-1}'=\lambda_y(\cS_{i+1}/\cS_{i-1})+B'_{i-1},
\end{gather*}
regarded as elements in \smash{$\QK_T^{\loc(\widehat{\mathbf{r}})}(Y)[y]$}.

By the projection formula, to prove~\eqref{eqn:whitY}, it suffices to prove the pushforward along $\pi$ of~\eqref{eqn:whitX} agrees with~\eqref{eqn:whitY}. We compare the two determinants by expanding along columns. Expanding along the column containing $B'_{i-1}$, we have that the determinant in~\eqref{eqn:whitY} is of the form
\begin{gather*}%\label{eqn:exprX}
 -B'_{i-1}\star C'+A'_{i-1}\star D'-E';
\end{gather*}
expanding along the two columns containing $B_{i-1}$ or $B_i$, we have that the determinant in~\eqref{eqn:whitX} is of the form
\begin{gather*}
 \begin{vmatrix}B_{i-1} & 0\\
 A_{i-1} & B_{i}\end{vmatrix}^\star\star0-\begin{vmatrix}
 B_{i-1} & 0\\
 1 & A_{i}
 \end{vmatrix}^\star\star C +\begin{vmatrix}
 B_{i-1} & 0\\
 0 & 1
 \end{vmatrix}^\star\star F\\
 \qquad{} +\begin{vmatrix}
 A_{i-1} & B_i\\1 & A_i
 \end{vmatrix}^\star\star D -\begin{vmatrix}
 A_{i-1} & B_i\\
 0 & 1
 \end{vmatrix}^\star\star E+\begin{vmatrix}
 1 & A_i\\
 0 & 1
 \end{vmatrix}^\star\star 0.
 \end{gather*}
Note that $C$, $D$, $E$, $F$ stay the same under the pushforward, and it is straightforward to check that
\[
C'=C,\qquad D'=D,\qquad E'=E.
\]
The rest follows from Lemma~\ref{lemma:main-push} below.
\end{proof}

\begin{Lemma}\label{lemma:main-push} The following hold:
 \begin{enumerate}[label=$(\alph*)$]\itemsep=0pt
 \item $\pi_* \begin{vmatrix}
 A_{i-1} & B_i\\
 0 & 1
 \end{vmatrix}^\star= 1$;
 \item $\pi_* \begin{vmatrix}
 B_{i-1} & 0\\
 0 & 1
 \end{vmatrix}^\star = 0$;
 \item $\pi_* \begin{vmatrix} B_{i-1} & 0\\ 1 & A_{i} \end{vmatrix}^\star = B_{i-1}'$;
 \item Assume that $r_i - r_{i-1} = 1$. Then
 \[
 \pi_* \begin{vmatrix} A_{i-1} & B_i\\1 & A_i
 \end{vmatrix}^\star = A_{i-1}'.
 \]
 \end{enumerate}
 \end{Lemma}
 \begin{proof}
 Note that $X$ may be realized as the Grassmann bundle $\mathbb{G}(r_i-r_{i-1}, \cS_{i+1}/\cS_{i-1})$
 over~$Y$, with tautological sequence $0 \to \cS_i/\cS_{i-1} \to \cS_{i+1}/\cS_{i-1} \to \cS_{i+1}/\cS_{i} \to 0$.
 It follows from Proposition~\ref{prop:relativeBWB} that
 \begin{equation}\label{eqn:pushforwards-lambda}
 \pi_*(\lambda_y(\cS_{i+1}/\cS_{i}))=\sum_{j=0}^{r_{i+1}-r_i}y^j\wedge^j(\cS_{i+1}/\cS_{i-1}), \qquad\pi_* (\lambda_y(\cS_i/\cS_{i-1})) = 1.
 \end{equation}

 For (a), (b), note that \smash{$A_{i-1}, B_{i-1}\in \QK_T^{\loc(\mathbf{\hat r})}(X)$}, so we may use~\eqref{E:Kato-loc-hom}, and it follows that
 \begin{gather*}%\label{eqn:pushforwards}
 \pi_*B_{i-1}=0,\qquad \pi_*A_{i-1} = 1.
 \end{gather*}
 Note that by Proposition~\ref{prop:MNS} and Theorem~\ref{thm:kato-morphism}, we have
 \begin{gather}
 \det\cS_{j}\star\det(\cS_{j+1}/\cS_{j})=(1-Q_{j})\det\cS_{j+1}\qquad \text{for $0\leq j\leq k$, in $\QK_T(X)$},\label{eqn:relX}\\
 \det\cS_{i-1}\star\det(\cS_{i+1}/\cS_{i-1})=(1-Q_{i-1})\det\cS_{i+1}\qquad \text{in $\QK_T(Y)$}.\label{eqn:relY}
\end{gather}
 To prove (c), we obtain from definition
 \begin{align}
 \begin{vmatrix}
 B_{i-1} & 0\\
 1 & A_{i}
 \end{vmatrix}^\star
 &{}= B_{i-1}A_i=B_{i-1} \star (\lambda_y(\cS_{i+1}/\cS_i)+B_i)\nonumber\\
 &{}= B_{i-1}\star \lambda_y(\cS_{i+1}/\cS_i)+ B_{i-1} \star B_i .\qquad\label{eqn:3-expand}
 \end{align}
 The element $B_{i}$ cannot be pushed forward, as it contains $1-Q_i$ in the denominator. However,
 we use~\eqref{eqn:relX} to calculate
 \begin{align*}
 B_{i-1} \star B_i &{}= y^{r_{i+1}-r_{i-1}} \frac{Q_{i-1} Q_i}{(1-Q_{i-1})(1-Q_i)} \det (\cS_{i+1}/\cS_i) \star \det (\cS_{i}/\cS_{i-1}) \\
 &{}= y^{r_{i+1}-r_{i-1}}Q_{i-1}Q_i\frac{\det\cS_{i+1}}{\det\cS_{i-1}} ,
 \end{align*}
 where the inverse is calculated in the quantum K ring of $X$. By~\eqref{eqn:relX} again,
 \[ \frac{\det\cS_{i+1}}{\det\cS_{i-1}} = \frac{\det \cS_{i+1} \star\det \C^n/\cS_{i-1}}{(1-Q_{i-1})\det \C^n}\qquad \text{in $\QK_T^{\loc(\mathbf{\hat r})}(X)$}, \]
 and its pushforward is
\begin{gather}\label{eqn:3}
 \frac{\det\cS_{i+1}}{\det\cS_{i-1}}\in \QK_T^{\loc(\mathbf{\hat r})}(Y).
\end{gather}
Note that by~\eqref{eqn:relY},
expression \eqref{eqn:3} is equal to
\begin{gather*}
 \frac{\det\left(\cS_{i+1}/\cS_{i-1}\right)}{1-Q_{i-1}}\qquad \text{in} \ \QK_T^{\loc(\mathbf{\hat r})}(Y).
\end{gather*}
Using \eqref{eqn:pushforwards-lambda} and \eqref{eqn:3-expand}, the projection formula, and Theorem~\ref{thm:kato-morphism}, it follows that
\begin{gather*}
 \pi_*\begin{vmatrix}
 B_{i-1} & 0\\
 1 & A_{i}
 \end{vmatrix}^\star=\pi_*\left(B_{i-1}\star B_i\right)=y^{r_{i+1}-r_{i-1}}\frac{Q_{i-1}}{1-Q_{i-1}}\det\left(\cS_{i+1}/\cS_{i-1}\right)=B'_{i-1}.
\end{gather*}
 For (d), we calculate
 \[ \begin{vmatrix}
 A_{i-1} & B_i\\1 & A_i
 \end{vmatrix}^\star
 = \begin{vmatrix}
 \lambda_y(\cS_{i}/\cS_{i-1}) + B_{i-1} & B_i\\1 & A_i
 \end{vmatrix}^\star
 = A_i \star \lambda_y(\cS_{i}/\cS_{i-1}) + A_i \star B_{i-1} - B_i. \]
 From (c), $\pi_* (A_i \star B_{i-1}) = B_{i-1}'$, therefore it suffices to show that
 $A_i \star \lambda_y(\cS_{i}/\cS_{i-1}) - B_i$ may be pushed forward, and that
 \begin{gather}\label{eqn:nts}
 \pi_*\left(A_i \star \lambda_y(\cS_{i}/\cS_{i-1}) - B_i\right) = \lambda_y(\cS_{i+1}/\cS_{i-1}) .
 \end{gather}
 The hypothesis $r_i - r_{i-1} =1$ implies that $\cS_i/\cS_{i-1}$ is a line bundle, and that
 \begin{gather*}%\label{E:int1}
 A_i \star \lambda_y(\cS_{i}/\cS_{i-1}) - B_i\\
 \qquad{} = \lambda_y(\cS_{i+1}/\cS_{i}) \star \lambda_y(\cS_{i}/\cS_{i-1})
 + y^{r_{i+1}-r_{i-1}} \frac{Q_i}{1-Q_i} \det (\cS_{i+1}/\cS_{i}) \star \det (\cS_{i}/\cS_{i-1}) .
 \end{gather*}
 By \eqref{eqn:pushforwards-lambda}, we have
 \[
 \pi_*(\lambda_y(\cS_{i+1}/\cS_{i}))=\lambda_y(\cS_{i+1}/\cS_{i-1})-y^{r_i-r_{i-1}}\det(\cS_{i+1}/\cS_{i-1}), \qquad\pi_*(\lambda_y(\cS_{i}/\cS_{i-1}))=1.
 \]
 By \eqref{eqn:relX}, we have
 \[
 \frac{Q_i}{1-Q_i} \det (\cS_{i+1}/\cS_{i}) \star \det (\cS_{i}/\cS_{i-1}) = Q_i(1-Q_{i-1})\frac{\det\cS_{i+1}}{\det \cS_{i-1}}.
 \]
 As in the proof of (c), this can be pushed forward and its pushforward is $\det(\cS_{i+1}/\cS_{i-1})$. Putting these together, we have established \eqref{eqn:nts}.
\end{proof}

Recall that $Y^{(j)}=\bigl(Y^{(j)}_1,\dots,
{ Y^{(j)}_{r_{j+1}-r_j}}\bigr)$, $0\leq j\leq k$, are formal variables, $e_\ell$ denotes the $\ell$-th elementary symmetric polynomial, and $T_1,\dots, T_n \in \K_T(\pt)$ are given by
the decomposition of~$\C^n$ into one dimensional $T$-modules, that is, $\wedge^\ell(\C^n)=e_\ell(T_1,\dots,T_n)$.

\begin{Theorem}\label{thm:main}
 The ring $\QK_T(\Fl(\bfr,n))$ is isomorphic to $R\dbb{Q}/J_Q$, where
 \[
 R=\K_T(\pt)\bigl[e_1\bigl(Y^{(j)}\bigr),\dots,e_{r_{j+1}-r_j}\bigl(Y^{(j)}\bigr),\, 0\leq j\leq k\bigr],
 \]
 and $J_Q\subset R[\![Q]\!]=R[\![Q_1,\dots,Q_k]\!]$
is the ideal generated by the coefficients of $y$ in
\begin{align}\label{eqn:whit-short-gen}
 \prod_{\ell=1}^{n}(1+yT_\ell)-
 &\begin{vmatrix}
 A_0 & B_1 & & & \\
 1 & A_1 & B_2 & & \\ & \ddots & \ddots & \ddots & \\ & & 1 & A_{k-1} & B_{k}\\ & & & 1 & A_{k}
 \end{vmatrix},
 \end{align}
 where
 \[
 A_j=\prod_{\ell=1}^{r_{j+1}-r_j}\bigl(1+y Y^{(j)}_\ell\bigr)+B_j,\qquad B_j=y^{{r_{j+1}-r_j}}\frac{Q_j}{1-Q_j}\prod_{\ell=1}^{r_{j+1}-r_j}Y^{(j)}_\ell,
 \]
 with the convention that $Q_0=0$.

 More precisely,
% for $\mathbf{r}=(1,2,\dots,n-1)$,
there exists a $\K_T(\pt)[\![Q]\!]$-algebra isomorphism
\[
\Psi\colon\ R[\![Q]\!]/J_Q \to \QK_T(\Fl(r_1,\dots,r_k))
\]
that sends \smash{$e_\ell\bigl(Y^{(j)}\bigr)$} to $\wedge^\ell\br{\cS_{j+1}/\cS_j}$ for $j=0,\dots,k$ and $\ell=1,\dots,r_{j+1}-r_{j}$.
\end{Theorem}

\begin{proof}
It follows from Theorem~\ref{thm:whit-short} that $\Psi$ is a well-defined ring homomorphism.
To prove there are no other relations, we use~\cite[Theorem~4.1]{GMSXZZ:Nakayama}, which
states that a complete set of relations in the quantum (equivariant) K ring is
obtained by quantizing any complete set of relations in the ordinary (equivariant) K ring.
Therefore, we need to show that when one specializes each~$Q_i$ to~$0$, the resulting
ring is a presentation of $\K_T(\Fl(r_1,\dots,r_k))$.
% Geometrically,
The relations obtained this way
are the `Borel-type relations' of the $\lambda_y$ classes
\begin{equation}\label{E:pToda} \lambda_y(\cS_1) \cdot \lambda_y(\cS_2/\cS_1)\cdot \dots \cdot \lambda_y(\C^n/\cS_k) = \lambda_y(\C^n) . \end{equation}
Note that the relations~\eqref{E:pToda} can be obtained from the Whitney relations
\[ \lambda_y(\cS_i) \cdot \lambda(\cS_{i+1}/\cS_i) = \lambda_y(\cS_{i+1}) , \]
by eliminating the classes $\lambda_y(\cS_i)$ for $2 \le i \le k$.
(The quantization of this statement is done in the next section.)
Finally, it is known that
the Whitney relations form a full set of relations
in $\K_T(\Fl(r_1,\dots,r_k))$. This is essentially done by Lascoux~\cite[Section~7]{lascoux:anneau},
and we refer to~\cite[Proposition~5.1]{GMSXZZ:Nakayama} for a complete proof.
\end{proof}
